\documentclass[conference]{IEEEtran}
\usepackage{iftex}
\ifXeTeX
\usepackage{fontspec}
\setmainfont{Times New Roman}
\setsansfont{Arial}
\setmonofont{Courier New}
\else
\usepackage[T1]{fontenc}
\usepackage{lmodern}
\fi
\usepackage{amsmath,amssymb,array,url,graphicx,etoolbox}
\IfFileExists{booktabs.sty}{\usepackage{booktabs}}{\newcommand{\toprule}{\hline}\newcommand{\midrule}{\hline}\newcommand{\bottomrule}{\hline}}
\patchcmd{\abstract}{---}{: }{}{}
\patchcmd{\IEEEkeywords}{---}{: }{}{}
\title{LLM-Mediated Natural Language Controls for Short-Video Recommender Systems: An Offline Study of Controllability}
\author{\IEEEauthorblockN{Vibhor Malik}\IEEEauthorblockA{M.Sc. in Applied Computing\\British Columbia Institute of Technology\\Supervisor: Borna Noureddin\\COMP 9500 Research Project\\20 October 2026}}
\begin{document}
\bstctlcite{controlNoAuthorDash}
\maketitle
\thispagestyle{plain}\pagestyle{plain}
\begin{abstract}
This offline study asks whether natural-language category controls can change a short-video feed and how they affect next-item ranking quality. It separates language parsing from a deterministic control policy. In Experiment 1, target-category exposure rises from 11.6\% to 79.2\%, while hard exclusion reduces muted exposure to zero. An earlier exposed practice run scored 73/100 validated Llama commands. A separate, prospectively frozen Experiment 2 uses 100 fresh synthetic requests and model-generated reference commands. Two separately contextualized model annotators agree on all 100 initial commands. A later self-reported manual response table descriptively matches 99/100 AI-authored initial labels after documented rubric encoding; formal adjudication remains incomplete. Llama matches 76 final references, compared with 11 for a rule parser; 23 Llama outputs fail validation and one accepted output mismatches its reference. Induced control success is 20.145\% for Llama versus 3.815\% for item feedback, while NDCG@10 is lower at 0.02580 versus 0.03184. Conditional requests achieve no control successes. These results show stronger pooled category control in this fixed synthetic setting, with a cost to historical next-item ranking. Model-generated references cannot establish human understanding or satisfaction.
\end{abstract}
\begin{IEEEkeywords}
controllable recommendation, natural-language interfaces, sequential recommendation, offline evaluation, reproducibility
\end{IEEEkeywords}

\section{Introduction}
Item feedback communicates a reaction to one recommendation. A request such as ``show more of this category'' specifies a broader intervention, while a profile can carry a preference across requests. Explanations serve a different purpose: making the information behind recommendations visible. Comparing these mechanisms requires separating the language model's interpretation from the ranker's mechanical enforcement.

Mozilla's 2022 audit of YouTube controls reported that Dislike prevented 12\% and Not interested 11\% of unwanted recommendations relative to its no-feedback control~\cite{ricks2022button}. This historical platform evidence motivates investigation; it is neither a current platform measurement nor a baseline for the present implementation. The research question here is narrower: how reliably does a local language model compile category requests, and how do category exposure and next-item accuracy change when the commands are applied to a frozen recommender?

This paper studies category-request parsing, deterministic control, and changes in exposure and next-item ranking on a frozen recommender. It uses historical public interactions and synthetic instructions without recruiting participants.

\section{Related Work}
Language-based profiles are an established direction. Radlinski et al. outline transparent, scrutable profiles~\cite{radlinski2022profiles}; Wo\'{z}niak et al. investigate user-editable profiles~\cite{wozniak2025profiles}; Ramos et al. study edits and recommendation changes~\cite{ramos2024transparent}. These motivate explicit preferences. Sanner et al., Chen and Pu, and Zhou et al. provide related profile and critiquing approaches~\cite{sanner2023competitive,chen2012critiquing,zhou2024language}. Niu et al. examine conversational short-video control through interviews and a user study~\cite{niu2025chat}; Fabbri et al. describe an integrated transparent, controllable feed~\cite{fabbri2026feeding}.

Lu et al. align language models with recommendation instructions using a conversational recommender with a backbone such as SASRec~\cite{lu2024aligning}. Dai et al. study learned control tokens for reranking~\cite{dai2025token}. Here, the recommender is frozen and the LLM translates commands. SASRec supplies the architectural basis for the compact network~\cite{kang2018sasrec}. KuaiRec provides historical short-video feedback and categories~\cite{gao2022kuairec}.

Evaluation design is central. Ji et al. identify leakage risks when recommendation splits admit future interactions~\cite{ji2023leakage}. A category derived from a held-out item is still outcome-derived information even if the item's identifier is hidden. The revised pipeline prevents that path. Zhang et al.'s vision paper motivates preference-alignment measures beyond historical ranking accuracy~\cite{zhang2026context}; it does not validate this project's numerical satisfaction bands. Compliance, predictive relevance and human preference remain distinct constructs.

\section{Method}
\subsection{Data and information boundaries}
The processed cohort contains 865 eligible users, 2,593 candidate items and 31 opaque official category identifiers. The source KuaiRec matrix contains 1,411 users, 3,327 items and 4,676,570 interactions~\cite{gao2022kuairec}. A watch ratio above 2.0 is a derived positive-engagement proxy, not an observed like. Official documentation defines watch ratio as play duration divided by video duration~\cite{kuairecdocs}.

Preprocessing excludes 181,992 records without timestamps rather than inventing dates. Shared chronological boundaries divide source events into nominal 70\% training, 15\% request history and 15\% evaluation windows, before positive-event filtering and eligibility. Users need at least three training, two request-history and one evaluation item. Candidates must have a positive event before the training cutoff; 43,708 post-cutoff positive cold-item events are removed before user filtering. The retained histories contain 127,173 training positives, 11,712 request-history positives and 4,662 test positives. These counts describe different stages and must not be interpreted as one common-denominator retention fraction.

Training uses only training sequences. Scoring receives training plus request history, truncated to its maximum context. Requests may access that history and static categories, never final behavioral labels. The first eligible final item supplies the relevance label only after assignment and ranking. The frozen full-catalogue policy permits seen items, although the retained cohort has no repeated training items or test/history overlap; repeated targets therefore do not justify that policy. An unseen-only candidate design would be a separate experiment.

Categories remain identifiers such as \texttt{category\_12}; the study does not infer topic names or test understanding of natural video subjects. Synthetic requests are deterministic experimental assignments, not observations of what users wanted. Historical support does not make them independent human preferences or establish a causal response to an intervention.

\subsection{Recommender and matched configurations}
The compact SASRec-style network uses item and position embeddings, two self-attention blocks, one attention head, hidden dimension 32, maximum history 50 and dropout 0.2. It uses Adam with learning rate 0.001, batch size 64 and gradient clipping at 5.0, for 20 fixed epochs in seeds 42, 43 and 44. Negative sampling excludes training-positive items without consulting later positives. Checkpoints remain fixed across control conditions. The initial 256-user, three-epoch run measured execution cost before control-outcome inspection; the main run's fixed larger scope is retained in its provenance. No current run retrains or tunes the backbone.

Five configurations use the same catalogue and frozen score vector: A is simulated positive or negative feedback on one historical exemplar; B adds display-only explanations to A; C compiles a text command; D compiles one explicit profile preference; E combines text, profile and explanations without double-counting a repeated command. B must equal A in rank order. Equivalent C, D and E commands deliberately share policy. These configurations do not identify every mechanism's independent contribution: a category command communicates more than a single-item reaction, and explanations have no offline ranking intervention.

\begin{table}[t]\centering\small
\caption{Matched mechanisms and the comparison each supports}
\begin{tabular}{p{0.4cm}p{2.6cm}p{3.9cm}}\toprule
ID&Mechanism&Interpretation\\\midrule
A&One historical item reaction&Explicit simulated reference policy\\
B&A with explanation display&Ranking invariance, not a preference effect\\
C&Parsed free-text command&Translation followed by enforcement\\
D&One profile preference&Equivalent-intent consistency within session\\
E&Text, profile and explanation&Merge consistency without double counting\\\bottomrule
\end{tabular}\end{table}

For each user, category support comes from training and request history. A hash selects supported boost and mute targets; the latest matching historical item supplies A's exemplar. The free-text and profile channels express the same category operation. Their canonical wording was used during development, so that task is a consistency check rather than held-out language generalization.

\subsection{Parsing, enforcement and explanations}
The actual language backend is local Llama 3.1 8B through Ollama, with constrained structured output and application validation. A transparent rule parser is a separate comparator, never a substitute for failed inference. For E1, a command contains boost, mute, reset or clarify, a category and strength. Unknown categories, unsupported operations and malformed outputs remain errors. In end-to-end ranking, an error applies no new control and stays in the denominator.

Scores are normalized by each user's catalogue minimum and maximum, with tied scores mapped to zero. The default category boost adds 0.25; a hard mute removes matching items without padding forbidden items back into an underfilled list. Item identifiers break ties deterministically. Profile state is session scoped, and explicit commands take precedence over an equivalent stored preference. The profile does not implement the proposal's rich, approximately 200-word human-edited representation.

Explanations select one to three verbatim statements from verified item, category and historical-association facts. This is grounded statement selection rather than unrestricted generation. Explanations do not change B's ranking.

The cache key binds verified facts to model provenance. Ten baseline items have cached model selections; other cards show a labeled verified template. The template is not counted as model inference.

Stateful interface checks are also separate from offline evidence. The API bounds inputs, isolates sessions/users and exposes backend/data status. Functional workflows cover setting preferences, changing users and resetting state. These checks can establish that the application applies a command consistently; they cannot determine whether a person understands a category identifier or prefers the resulting feed. A persistent profile across application restarts and a licensed video-playback experience are outside the implemented prototype.

\section{Evaluation}
\subsection{Metrics and statistical units}
At $K=10$, target-category proportion TCP@10 is the fraction of slots carrying a boosted target; TCER@10 is the same arithmetic for an unwanted target, where lower is better. The harness separately reports fill rate, conditional exposure for nonempty lists and filled compliance. Zero exposure after an empty output must not be mistaken for successful control. A positive presence check does not guarantee improved exposure over baseline.

Relevance uses one historical next-item label. HR@1 is one when it ranks first; NDCG@10 is $1/\log_2(r+1)$ when it appears at rank $r\leq10$, otherwise zero. These are next-item recovery proxies, not utility over all candidates or human satisfaction. A steering request can legitimately disagree with the historically observed next item, which predates the intervention.

E1 averages paired outcomes across requests and three seeds within each user. Thus 865 users remain 865 analysis units, rather than 5,190 independent observations. Its six primary contrasts are C, D and E versus A for TCP and TCER, with Wilcoxon signed-rank tests, Pratt zeros, rounded ties and Benjamini--Hochberg adjustment~\cite{scipywilcoxon,scipyfdr}. User-bootstrap intervals use 5,000 resamples. Shared configurations, sparse nonzero differences, ties and dependence restrict interpretation. Neither a small $p$ value nor a narrow interval corrects synthetic-intent bias or a weak backbone.

\subsection{Language accuracy and non-inferiority}
The frozen E1 broader suite measures exact operation/category/strength agreement with its fixed synthetic rubric. It maps ``less'' and ``fewer'' to hard mute, treats conditional requests as unsupported and asks for clarification when a neutral ``show'' request lacks direction. Those conventions are not established human meanings. One development/test pair differs only by a full stop; shared templates further limit independence. The separate 120 canonical inputs were development wording, and the ten adversarial cases form a small finite stress set.

The original parser target is 85\%, with a 70\% degradation floor. The broader suite's descriptive 95\% Wilson interval assumes a binomial working model; it does not establish population-wide reliability. Operational and validation failures remain in the denominator. Request/seed reuse cannot enlarge the number of unique language observations. A later prompt change would create development evidence and require a fresh prospective bank, not repair the interpretation of an already observed test.

The parser report distinguishes three levels: a model response may complete, its JSON may satisfy the supplied schema, and the validated command may agree with the reference. These are different outcomes. Application grounding can reject a syntactically valid command, while a valid accepted command can still express the wrong operation. End-to-end exact agreement treats both cases as failures. Reporting only completed HTTP requests or only cases that pass validation would inflate reliability. Raw responses, accepted commands, rejected commands and clarification outcomes must therefore remain traceable per question.

The proposal's ``within 10\%'' non-inferiority phrase is ambiguous. The frozen implementation used absolute HR@1 margin 0.10; the later audit additionally reports 10\% of the paired A baseline as a post hoc sensitivity analysis. With baseline about 0.0025048, that relative margin is about 0.00025048. The simultaneous lower bound for C/D/E minus A is about $-0.0007707$: above $-0.10$, below $-0.00025048$. The absolute calculation passes numerically but could permit complete baseline loss; the relative reading is unsupported. Neither establishes useful quality preservation or retroactive prespecification.

\section{Experiment 1: Category Control and Ranking}
The executed cohort and checkpoints remain unchanged. Table~\ref{tab:e1} summarizes the matched control run. Target exposure changes are large, but policy construction explains their direction: category-wide boosting directly favors TCP against item-only feedback, and hard exclusion guarantees zero muted exposure after correct parsing. C, D and E have identical rankings on equivalent canonical commands; the separately executed rule comparator gives the same recorded rankings and aggregate metrics. This attributes the canonical-task gain to shared enforcement, not an observed LLM advantage.

\begin{table}[t]\centering\small
\caption{Frozen E1 control results, user means}\label{tab:e1}
\begin{tabular}{lrr}\toprule
Metric & A/B & C/D/E\\\midrule
Boost TCP@10 (\%) &11.615&79.222\\
Mute TCER@10 (\%) &6.112&0.000\\
Fill rate (\%) &100.000&100.000\\
Boost NDCG@10 &0.03012&0.01637\\
Mute NDCG@10 &0.03370&0.03526\\\bottomrule
\end{tabular}\end{table}

For C versus A, the mean paired boost TCP difference is 67.607 percentage points, with 95\% user-bootstrap interval [65.788, 69.438]; the TCER difference is $-6.112$ points, interval [$-7.056,-5.241$]. D and E yield the same contrasts. The BH-adjusted Wilcoxon values are approximately $1.35\times10^{-142}$ for boosts and $1.05\times10^{-65}$ for mutes. The proposal's directional control criterion is numerically supported under these conventions; this is not evidence that the LLM improves recommendation accuracy or understands general human intent.

Quality moves differently. Boost NDCG@10 decreases 45.7\%. The post hoc paired boost C--A difference is $-0.01375$, with user-bootstrap interval [$-0.02086,-0.00655$], and 108 users have nonzero differences. This secondary analysis was added after inspection and remains outside the unchanged primary BH family. Muted requests' NDCG increases, so one pooled metric would conceal distinct effects. The historical next-item label may disagree with a newly assigned steering instruction; the result still rules out an unqualified ``no quality cost'' claim.

The backbone is weak: per-seed HR@1 is 0.000\%, 0.347\% and 0.462\%, against about 7.977\% for request-window popularity. The latter uses pooled recent pre-test histories; the trained network uses each user's recent history for scoring but its weights were fitted on the earlier training window. Both avoid final labels, but this estimator difference prevents a controlled model-class superiority claim. Between 73.9\% and 96.2\% of users share the most-common first prediction within a seed. Such concentration and near-zero accuracy make non-inferiority especially uninformative.

\begin{table}[t]\centering\small
\caption{Historical uncontrolled backbone diagnosis}
\setlength{\tabcolsep}{3pt}
\begin{tabular}{lrrr}\toprule
Model&HR@1 (\%)&HR@10 (\%)&NDCG@10\\\midrule
Uniform, expected&0.039&0.386&0.00175\\
Training popularity&0.116&0.578&0.00298\\
Recent popularity&7.977&16.532&0.11866\\
SASRec seed 42&0.000&4.509&0.01642\\
SASRec seed 43&0.347&10.058&0.03627\\
SASRec seed 44&0.462&11.098&0.04584\\\bottomrule
\end{tabular}\end{table}

Llama parses all 120/120 canonical inputs, 79/100 broader synthetic requests and 2/10 adversarial requests. The broader-set Wilson interval is [70.02\%, 85.83\%]: the 70\% floor is cleared, while the 85\% point target is missed. The rule comparator obtains 81/100, not proof of general superiority from this finite set. Seven Llama cases are grounding rejections, recorded as errors rather than discarded; correct structured syntax and intended meaning are separate checks. All ten generated explanation cards equal the template, establishing no observed explanation advantage in this limited sample.

\begin{table}[t]\centering\small
\caption{Frozen E1 broader synthetic language results}
\begin{tabular}{lrr}\toprule
Case group&Rule correct&Llama correct\\\midrule
Positive&26/30&30/30\\
Negative&24/30&30/30\\
Reset&3/10&10/10\\
Ambiguous/unsupported&18/20&7/20\\
Adversarial&10/10&2/10\\\midrule
Total&81/100&79/100\\
Validation rejections&0&7\\\bottomrule
\end{tabular}\end{table}

The supported positive, negative and reset cases are all correct for Llama under the fixed E1 rubric. Difficulty concentrates in ambiguity and adversarial instructions. Fourteen broader-set cases yield valid but rubric-incorrect commands, and seven are rejected.

Historical verification reconstructs prepared sequences and recomputes all 25,950 actual-model ranking rows. A later audit reproduces all 25,950 top-ten lists but reports 2,009 full-ranking hash differences and 5,190 request-score hash differences. Their cause remains unresolved; missing original score vectors prevent quantifying numerical error. Dependency pins alone do not guarantee cross-platform equality. These records are historical Linux checks, not a newly completed E1 replay on Windows or Colab.

\section{Exposed Practice Evaluation}
The completed practice run at commit \texttt{c6ebda0} used 100 already exposed worked questions and model-generated answers,25 in each of four intent classes. It is preserved unchanged and was not rerun. Primary validated whole-command agreement remains 73/100 for Llama and 18/100 for rules. Llama's 24 errors comprise 23 conditional-grounding rejections and one floor rejection; raw pre-validation agreement 95/100 is supplementary, never a replacement score. Conditional validated agreement is 2/25, compared with 25/25 compound,22/25 floor and 24/25 graded cases.

The practice run retains 3,459 assigned requests,41,508 ranking rows and 2,595 full score vectors. User-averaged induced success is 23.81\% for Llama,8.40\% for rules and 3.92\% for item feedback. Llama NDCG@10 is 0.02461 versus item feedback 0.03128, so the observed control benefit accompanies a predictive quality cost. Recency reaches 0.11866 NDCG@10. A separate standard-library recount reproduces all row metrics and pooled means without discrepancies; that earlier audit did not independently reimplement intervals or signed-rank tests. Exposed questions cannot support an unseen-test claim or enlarge the 100 unique language observations through request/seed reuse.

\section{Evaluation with Model-Generated References}
\subsection{Design and provenance}
The fresh synthetic-request study followed a protocol frozen before cases were authored.

The protocol was frozen at 20:44:17 UTC on 14 October 2026, before fresh wording. It fixes the research question, sampling, prompts, rubric, failure handling, metrics and disagreement procedure. A seeded 100-slot plan balances compound, graded, nonzero-floor and same-item conditional requests (25 each), covering all 31 target identifiers. The grade mix is nine slight boosts, eight strong reductions and eight exclusions. A separate model author supplied three wording candidates per slot. First passing candidates were selected in fixed order without parser predictions.

Screening excludes exact normalized duplicates and close lexical matches to the exposed 100 questions and 48 archived development texts. The prespecified flags are normalized SequenceMatcher similarity at least 0.82 or token-bigram Jaccard at least 0.55, after category placeholders and punctuation normalization. All candidates and exclusions remain saved; fresh within-bank near-paraphrase flags are reported, not selected on outcomes. These heuristics do not guarantee semantic or pretraining novelty.

Two separately contextualized model annotators receive independently shuffled question sheets and the fixed semantic rubric, without author plans, implementation, predictions or each other's labels. Their original labels are locked before comparison; agreement supplies the final reference, while an additional model context would adjudicate a disagreement without predictions. Incompatible or unresolved references stop inference. Both passes use the same reported GPT-6.1-sol model family, so shared errors remain possible. These are model-generated references, not authenticated human labels.

The first reference-custody attempt aborted before comparison or inference because of a receipt-field mismatch. Its originals remain archived. A separately frozen second attempt reused the same bank and semantic protocol, with two new model annotation passes. The custody record preserves this chronology without silently repairing the failed receipts.

\subsection{Execution and metrics}
All final cases, both labels, adjudication, sources, prompts, validators, ranking inputs, model digest, seeds and metrics receive a full pre-inference hash lock. The Llama runner receives only opaque IDs, exact requests and the 31 public category IDs. Trusted worker code applies file/network audit restrictions and verifies denied reads of the protected label files; no reference contents enter Ollama's messages. Allowed networking is loopback 11434 only. This is application separation for trusted code on a shared Windows host, not an OS security sandbox against a compromised administrator.

The unchanged local \texttt{llama3.1:8b} / \texttt{Q4\_K\_M} model has SHA-256 digest \texttt{46e0c10c039e0191}\allowbreak\texttt{19339687c3c1757c}\allowbreak\texttt{c81b9da49709a3b3}\allowbreak\texttt{924863ba87ca666e}, using Ollama 0.34.0. Fixed options are temperature 0, seed 42, context 4096, output limit 768, six threads and 600-second timeout. Rules and Llama each attempt every exact request once. Durable pre-call journals and terminal records prevent silent reruns; any interruption aborts the version. Failures apply no new ranking control and remain in parser/ranking denominators. All 200 terminal predictions are committed before protected scoring joins references. No prompt, validator, band or threshold is tuned from practice or new outcomes.

The same 865-user,2,593-item data and three fixed SASRec checkpoints (42/43/44) supply matched item-feedback, rules, Llama and recency conditions. Eligibility depends on pre-test histories and final model-reference categories; SHA-256 assignment seed 2026 selects one eligible request per user/class. Next-item labels are accessed only for relevance after rankings. No retraining occurs. Let $D(S,L)$ be top-ten exposure to category set $S$ in list $L$, with each matching rank weighted by $1/\log_2(r+1)$ and normalized by the sum of those ten rank weights. Signed movement $M$ compares $D$ for the controlled list with the uncontrolled ranking, reversing direction for reductions and exclusions. Compound requests take the minimum directive movement; conditional requests subtract any added target exposure outside the condition; a floor violation sets movement to zero. Induced success requires ten returned items, the fixed class-specific category-count band and $M>10^{-12}$. These numerical bands are conventions rather than human satisfaction judgments.

Users are analysis units, averaging three seeds within each request/class and eligible classes equally within user. Paired 5,000-resample percentile bootstrap intervals use seed 2026. Eleven one-sided Pratt Wilcoxon comparisons use 12-decimal rounding, asymptotic ties and joint BH correction: four classes plus pooled Llama--A and Llama--rules success, and pooled Llama--A movement. H1\_AI needs both positive significant pooled success contrasts; H2\_AI needs positive significant movement. HR@1 and NDCG@10 are reported together with control outcomes as descriptive H3\_AI, without a noninferiority claim. Complete condition/seed masks preserve unfavorable outcomes, infeasibility, missing opportunity and exclusions.

\subsection{Results}
Initial whole-command agreement between model annotation passes is 100/100, Cohen kappa 1.000000; 0 disputes required the prescribed model adjudication. Directive and field agreement, original labels, rationales and decisions remain in the supplement. This measures consistency between model contexts, not factual correctness.

\begin{table}[t]\centering\small
\caption{Parser agreement with model-generated references}\label{tab:aiparser}
\begin{tabular}{lrrr}\toprule
Class & Rules & Llama & Rejected\\\midrule
Compound & 6/25 & 25/25 & 0\\
Graded & 5/25 & 24/25 & 0\\
Floor & 0/25 & 25/25 & 0\\
Conditional & 0/25 & 2/25 & 23\\
All & 11/100 & 76/100 & 23\\\bottomrule
\end{tabular}\end{table}

The run preserves 3,460 assigned user/class requests,41,520 ranking rows and 2,595 full score vectors. There are 0 eligibility exclusions. Across 10,380 request-seed assignments, 987 are structurally infeasible and 5,586 have no directional opportunity; these flags may overlap and remain in the main denominator. All 23 Llama validation rejections concern condition grounding. Parser errors remain in 100-case denominators and every assigned ranking comparison.

\begin{table}[t]\centering\small
\caption{Fresh synthetic-request ranking outcomes, paired user means}\label{tab:airank}
\begin{tabular}{lrrr}\toprule
Condition & Success(\%) & HR@1(\%) & NDCG@10\\\midrule
Item feedback & 3.815 & 0.241 & 0.03184\\
Rules & 3.854 & 0.231 & 0.03097\\
Llama & 20.145 & 0.337 & 0.02580\\
Recency & 5.202 & 7.977 & 0.11866\\
\bottomrule
\end{tabular}\end{table}

Prespecified pooled directional tests support H1\_AI and H2\_AI after joint Benjamini--Hochberg correction across 11 contrasts. H3\_AI reports HR@1 and NDCG@10 descriptively.

The mean user success rate is 20.145\% for Llama, 3.854\% for rules and 3.815\% for item feedback. Conditional requests have 0.000\% induced success in all four conditions, including Llama, which matches only 2/25 conditional reference commands after validation. The pooled gain comes from compound, graded and floor requests; it does not establish conditional control. NDCG@10 is 0.02580 for Llama versus 0.03184 for item feedback. Control and next-item ranking quality must be interpreted together.

The saved computational audit recalculates the parser and ranking results without discrepancies within its stated scope. Reused questions across users and seeds do not increase the 100 distinct language observations.

\section{Validity and Reproducibility}
Synthetic opaque categories restrict construct validity. A single historical exemplar conveys less information than a category instruction. Equivalent policy outputs do not establish equivalent usability, while hard exclusion's zero exposure is true by construction. The study neither observes real steering intentions nor measures human comprehension or satisfaction. KuaiRec's platform, period, positive-event threshold and filtered active cohort restrict external validity.

Shared templates and exposure limit language independence. The weak untuned backbone, sparse paired quality differences and globally pooled recency comparator complicate quality interpretation. The original absolute and relative non-inferiority readings remain visible. The exposed practice bank and computational audit answer different questions from an unseen test or human annotation. Older ranking and score hashes still need reconciliation, although top-ten rankings reproduced.

The artifact separates preprocessing, training, controls, evaluation and a React/FastAPI demonstration. Real-data versus fixture status and rule versus local-model status are visible. Metadata cards are not playable video assets. State is session scoped; malformed commands remain errors. Prior unit/API/browser checks support software behavior, not the research hypotheses or present full-platform acceptance. The Windows practice and prospective model-reference runs are verified; full E1 replay and Colab execution remain unverified.

The separate local evidence archive retains the literal prompt and schema, runtime digest, input hashes, model-generated references, raw responses, parser records, score vectors, ranking rows, terminal logs and audit results. It records an interrupted wrong-model attempt whose outputs were excluded. The saved run used Ollama with the archived Llama digest, Python 3.14.0, NumPy 2.3.5, SciPy 1.16.3 and PyTorch 2.9.1 CPU. Saved records permit computational replay without new model calls; a different inference runtime may change outputs. The public repository holds compact evidence and does not redistribute model weights.

\noindent\textbf{Subsequent response table (20 October 2026).} Vibhor Singh Malik reports manually completing all 100 questions in the original v2 student worksheet without seeing reference answers. The original CSV has 100 unique expected IDs and 125 parseable directives. After submission, he approved applying the existing rubric convention to 99 unqualified boost or reduce strengths and eight exclusion strengths; 18 explicitly graded strengths were kept. The raw submission and separately frozen encoding remain distinct. The encoded commands all pass schema validation. The recovered original v2 bank matches the digest in its prior commitment and package manifests. A separate descriptive comparison finds 99/100 exact whole-command matches with the bank's AI-authored initial labels, Cohen's kappa 0.9899, and 124/125 matching strength fields. The single differing strength was explicitly entered, not one of the later rubric encodings; neither original label was changed. This is not inter-human agreement. Exact annotation times and the formal independent-pass receipt were not recorded, and historical development-summary exposure reports still need an eligibility account. No blind adjudication or original-route parser and ranking evaluation follows from this comparison. The completed 14 October model-reference results above are unchanged.

\section{Conclusion}
Experiment 1 demonstrates category enforcement and a predictive quality cost, with no LLM-specific ranking advantage on canonical requests. The exposed practice result remains a separate 73/100 validated diagnostic. In the fresh model-reference evaluation, Llama agrees on 76/100 whole commands versus 11/100 for rules, with 23 validation rejections and one accepted mismatch. Its mean user control success is 20.145\% versus 3.815\% for item feedback, while NDCG@10 falls to 0.02580 versus 0.03184. The prespecified pooled H1\_AI and H2\_AI directional tests pass, but conditional requests achieve no control successes. References generated by the same model family do not establish independent human truth, naturalness, satisfaction or performance across LLMs. The later student response table matches 99/100 AI-authored initial labels descriptively, but formal pass documentation and adjudication remain incomplete. Human experience remains a separately designed extension.

\bibliographystyle{IEEEtran}
\bibliography{references}
\end{document}
